% GENERATED: edit thesis/ch4_paper3.tex or paper3/frontmatter.tex.
\documentclass[11pt]{article}
\usepackage[margin=1.1in]{geometry}
\usepackage{amsmath,amssymb,amsthm,mathtools}
\usepackage{booktabs,threeparttable,graphicx,setspace,enumitem,longtable,array}
\usepackage{xcolor}
\usepackage{etoolbox}
\usepackage{placeins}
\AtBeginEnvironment{tabular}{\footnotesize}
\usepackage[authoryear,round]{natbib}
\usepackage[colorlinks=true,linkcolor=blue!50!black,citecolor=blue!50!black,urlcolor=blue!50!black]{hyperref}
\onehalfspacing
\emergencystretch=2em
\setcounter{tocdepth}{1}

\newtheorem{theorem}{Theorem}[section]
\newtheorem{proposition}[theorem]{Proposition}
\newtheorem{lemma}[theorem]{Lemma}
\newtheorem{corollary}[theorem]{Corollary}
\newtheorem{conjecture}[theorem]{Conjecture}
\theoremstyle{definition}
\newtheorem{definition}[theorem]{Definition}
\newtheorem{assumption}[theorem]{Assumption}
\theoremstyle{remark}
\newtheorem{remark}[theorem]{Remark}
\newtheorem{example}[theorem]{Example}
\newcommand{\E}{\mathbb{E}}
\newcommand{\R}{\mathbb{R}}
\newcommand{\KL}{D_{\mathrm{KL}}}
\newcommand{\lab}[1]{\textsc{[#1]}}

% status tags used in the new chapters
\newcommand{\tagbox}[2]{\fcolorbox{#1}{#1!7}{\textsf{\scriptsize #2}}}
\newcommand{\Proved}{\tagbox{green!45!black}{PROVED IN SOURCE}}
\newcommand{\ProvedIn}[1]{\tagbox{green!45!black}{PROVED IN #1}}
\newcommand{\Target}{\tagbox{blue!70!black}{TARGET: ACCEPTANCE CONDITIONS OPEN}}
\newcommand{\Conj}{\tagbox{orange!85!black}{CONJECTURE}}
\newcommand{\Refuted}{\tagbox{red!70!black}{REFUTED AS STATED}}
\newcommand{\Pending}{\tagbox{red!70!black}{PENDING DATA}}
\newcommand{\pendbox}[2]{\par\medskip\noindent\fcolorbox{red!70!black}{red!4}{\parbox{0.96\linewidth}{\small\textbf{\color{red!70!black}Pending --- #1.} #2}}\par\medskip}

% each chapter's references appear as a section at the end of that chapter
\renewcommand{\bibsection}{\section*{References}}

\begin{document}
\title{Certify, Don't Identify: Safety Without Identification in Adaptive Markets}
\author{Pratik Niroula}
\date{Working draft; unresolved targets and evidence are labelled}
\maketitle
\begin{abstract}
We study operational safety in markets with adaptive trading agents when identities are not recoverable. Observable response taps identify a merged memory measure, while safety depends on a low-dimensional functional. Under reset experiments, we give a coefficient-energy upper bound, a best-approximation obstruction and a long-memory impossibility result. A uniform matching lower bound over the proposed regular coefficient class fails; we give a counterexample and a weaker fixed-atom bound. We extend safety certification to signed, nonlinear and time-varying agents and derive a certificate lifetime under bounded weight drift. A positive-system comparison yields a sharper multi-asset bound, with registered numerical validation and its remaining formal boundary stated explicitly. A threshold-game counterexample shows that convex costs alone do not guarantee a strategic equilibrium. Formal proofs cover the finite-horizon safety induction and deterministic lifetime and interval arithmetic. Prospective empirical tests remain pending.
\end{abstract}
\section{Introduction}
The identification companion paper \citep{p3:paper1} studies what observable prices and aggregate flow reveal about adaptive trading cohorts. Its merged-measure equivalence separates observable aggregate dynamics from unidentifiable cohort decompositions. Moment-matched alternatives make cohort separation costly, while class-tagged disclosure changes the observation contract. These identification limits motivate asking which operational questions can be answered without recovering the underlying agents.

The safety companion paper \citep{p3:paper2} studies transient loss and an uncertainty-aware pre-trade guard. In its positive scalar economy, the summed impulse response determines the worst-case loss under bounded shocks. Certification can therefore be substantially cheaper than identifying the agents, conditional on a memory bound. Its observational forecasting evidence does not establish a causal liquidation mechanism or a universal real-market guarantee.

The present paper develops the common functional framework and extensions for changing agents. The theory below is generated from the thesis's authoritative Paper 3 source, with proofs collected in the appendix. Status boxes distinguish written results, open targets and pending empirical work. The general matching lower-bound target needs reformulation; the positive matrix result still has an explicit formalization boundary. The strategic counterexample concerns the abstract payoff assumptions as stated, not every economically restricted trading game.

\section{Markets with adaptive agents}\label{sec:fw_model}

\begin{definition}[Cohort economy]\label{def:cohort}
A state $x_t\in\R$, the displacement of the price from its fundamental value, evolves as
\begin{equation}\label{eq:fw_plant}
x_{t+1}=a\,x_t+\kappa\,(u_t+w_{t+1}),\qquad u_t=\sum_{i=1}^K q_i\,\varphi_i(y_{i,t}),\qquad y_{i,t}=r_i\,y_{i,t-1}+(1-r_i)\,x_t,
\end{equation}
with persistence $a\in[0,1)$, impact $\kappa>0$, exogenous shocks $w_t$, and $K$ cohorts of agents with weights $q_i\ge0$, memory rates $r_i\in[0,1)$ and response rules $\varphi_i$. The economy is \emph{linear} when $\varphi_i(y)=y$ for all $i$.
\end{definition}

The exponential state $y_{i,t}$ is the reduced form of steady-state Kalman filtering and of regularized online gradient descent (Lemma~\textup{2.1 in Paper~1}); Section~\ref{sec:bl_robust} relaxes linearity and stationarity. In the linear economy the aggregate response is a linear filter of past states,
\[
u_t=\sum_{m\ge0}k_m\,x_{t-m},\qquad k_m=\int r^m\,dM(r),\qquad M=\sum_{i=1}^K q_i(1-r_i)\,\delta_{r_i},
\]
so the agents enter the observable dynamics only through the \emph{merged measure} $M$, a positive measure on $[0,1)$, and through its moments $k_m$. Cohorts that share a memory rate are indistinguishable, whatever their identities (Theorem~\textup{3.1 in Paper~1}). We write $\mathcal M(\bar r,B)$ for the positive measures on $[0,\bar r]$ of total mass at most $B$; $\bar r<1$ is a \emph{memory bound}.

\section{Questions are functionals of the merged measure}\label{sec:fw_questions}

Every question an outside observer can answer from the observable dynamics is a functional $F$ of $M$ (and of the known plant parameters $a,\kappa$). Three families matter.

\paragraph{Identification.} $F(M)=M$ itself, in a metric such as the Wasserstein distance $W_1$, or the atoms $(r_i,q_i)$ when $M$ has $K$ atoms. Paper~1 studies these.

\paragraph{Safety.} For bounded shocks $|w_t|\le\bar w$, the worst-case transient loss of a position $I$ is $|I|\,\bar w\,\|h\|_1$, where $h$ is the closed-loop impulse response, and a pre-trade guard on $\|h\|_1$ admits every safe action and no unsafe one (Lemma~\textup{6.1 in Paper~2}). In the linear economy with $a\ge0$ all $h_t$ are non-negative, so $\|h\|_1$ equals the summed response
\begin{equation}\label{eq:fw_S}
S(M)=\frac{\kappa}{1-a-\kappa\,\Phi(M)},\qquad \Phi(M)=\sum_{m\ge0}k_m=\int\frac{dM(r)}{1-r}=\sum_i q_i ,
\end{equation}
whenever $a+\kappa\Phi<1$ (Theorem~\textup{4.4 in Paper~2}). Safety therefore depends on the agents only through one linear functional of $M$, the total feedback weight $\Phi$.

\paragraph{Other operational questions.} The response at a given lag, $k_m=\int r^m dM$; the discounted feedback $\int (1-\beta r)^{-1}dM$; the peak of the impulse response. The first two are linear in $M$; the peak is not, and a guard built on it admits unsafe actions (Paper~2).

\section{Observation contracts and the cost of a question}\label{sec:fw_contracts}

\begin{assumption}[Reset contract]\label{ass:reset}
The observer can run $N$ independent reset experiments. Each returns the first $L$ aggregate response taps with independent Gaussian noise, $\hat k_m=k_m+\sigma\xi_m$, $\xi_m\sim N(0,1)$, $m=0,\dots,L-1$.
\end{assumption}

This is the controlled-probe contract with observed aggregate flow (contract $I_1$ of Paper~1). Paper~2 works with the closed-loop taps $h_t$ instead; for the summed response the two contracts give the same rate (Theorem~\textup{5.5 in Paper~2}). Correlated tap noise, passive data and tagged disclosure are other contracts, treated in Chapters~1 and~2.

\begin{definition}[Cost of certification]\label{def:ncert}
For a functional $F$, a class $\mathcal M$, accuracy $\epsilon>0$ and level $1-\alpha$, $N_{\rm cert}(\epsilon;F,\mathcal M)$ is the smallest $N$ for which some interval built from the data has half-width at most $\epsilon$ and covers $F(M)$ with probability at least $1-\alpha$ for every $M\in\mathcal M$.
\end{definition}

\section{The certification--identification dichotomy}\label{sec:fw_dich}

Consider a linear question $F_g(M)=\int g\,dM$, where $g(r)=\sum_{t\ge0}a_t r^t$ converges absolutely on $[0,\bar r]$. Two numbers describe how hard it is. Its \emph{tail at depth} $T$ and its \emph{coefficient energy} are
\[
\tau_g(T)=B\sup_{0\le r\le\bar r}\Bigl|\sum_{t\ge T}a_t r^t\Bigr|,\qquad A_g(T)=\sum_{t<T}a_t^2,
\]
and its \emph{truncation depth} is $T_g(\epsilon)=\min\{T\ge1:\tau_g(T)\le\epsilon/2\}$. We also write $E_T(g)=\min_{\deg p<T}\max_{[0,\bar r]}|g-p|$ for the error of the best polynomial approximation of degree below $T$.

\begin{theorem}[Certification of linear questions]\label{thm:dich}
Under Assumption~\ref{ass:reset}, with $0<\alpha<1/2$, $L\ge T_g(\epsilon)$, and $z=z_{1-\alpha/2}$:
\begin{enumerate}[label=(\roman*)]
\item \emph{Upper bound.} $N_{\rm cert}(\epsilon;F_g,\mathcal M(\bar r,B))\le\bigl\lceil4z^2\sigma^2A_g(T_g(\epsilon))\,\epsilon^{-2}\bigr\rceil$. \hfill\Proved
\item \emph{Truncation is necessary.} For every $T$ there are $M,M'\in\mathcal M(\bar r,B)$ whose first $T$ taps coincide and with $|F_g(M)-F_g(M')|\ge2B\,E_T(g)$. Hence any certificate that uses only the first $T$ taps has worst-case half-width at least $B\,E_T(g)$. \hfill\Proved
\item \emph{Long memory makes certification impossible.} Without a memory bound, that is on the positive measures on $[0,1)$ of mass at most $B$, if $\limsup_{r\uparrow1}|g(r)|\sqrt{1-r}=\infty$ then for every level $1-\alpha$ with $\alpha<1/4$ no finite $N$ certifies $F_g$ to any finite accuracy. \hfill\Proved
\item \emph{Matching lower bound.} For $g$ in a regular class (coefficients of eventually constant sign and at most polynomial growth), $N_{\rm cert}(\epsilon;F_g,\mathcal M(\bar r,B))\ge c\,\sigma^2A_g(c'\,T_g(\epsilon))\,\epsilon^{-2}$ for constants $c,c'>0$ depending only on $(\bar r,B,\alpha)$. \hfill\Refuted
\end{enumerate}
\end{theorem}

\paragraph{Proof.} See Appendix~\ref{app:paper3proof1}.


\paragraph{What the theorem says.} The upper bound for answering a linear question is controlled by its Taylor coefficients and their tail; finite-depth ambiguity is controlled separately by best polynomial approximation. If its power-series tail decays geometrically on $[0,\bar r]$, the truncation depth grows like $\log(1/\epsilon)$ and the cost is parametric up to that factor. If the question is singular at infinite memory, faster than $(1-r)^{-1/2}$, no memory bound means no certificate. The number of agents and the spacing of their memories do not enter.

\begin{corollary}[The rates of this thesis as cases of one picture]\label{cor:cases}
\begin{enumerate}[label=(\alph*)]
\item \emph{Safety with a memory bound.} For $\Phi$, $a_t\equiv1$, so $A_\Phi(T)=T$ and $\tau_\Phi(T)=B\bar r^T/(1-\bar r)$. Part (i) gives $N_{\rm cert}=O\bigl(\sigma^2\epsilon^{-2}\log(1/\epsilon)/\log(1/\bar r)\bigr)$, and $S=\kappa/(1-a-\kappa\Phi)$ is certified by mapping the interval through this monotone function. This is the rate of Theorem~\textup{5.5 in Paper~2}, which Theorem~\textup{6.5 in Paper~2} shows is optimal up to the logarithm for the guard. \hfill\ProvedIn{PAPER 2}
\item \emph{Safety without a memory bound.} $g(r)=(1-r)^{-1}$ satisfies the condition of (iii): no finite number of experiments certifies the risk, as Theorem~\textup{5.3 in Paper~2}(iii) shows for finite windows. \hfill\ProvedIn{PAPER 2}
\item \emph{A single response tap.} For $g(r)=r^m$, $A_g=1$ and the cost is exactly parametric. \hfill\Proved
\item \emph{Identifying $K$ cohorts.} Recovering the atoms is not a linear question. It requires moments up to order $2K-1$ to be resolved against moment-matched alternatives whose higher moments differ by $O((\Delta\eta)^{2K-1})$, which costs $(\Delta\eta)^{-(4K-2)}$ experiments (Theorems~\textup{6.4 in Paper~1} and~\textup{6.5 in Paper~1}). \hfill\ProvedIn{PAPER 1}
\item \emph{Identifying a general memory measure.} Recovering $M$ in $W_1$ means certifying $F_g$ uniformly over all $1$-Lipschitz $g$. Their best approximations by degree-$T$ polynomials improve only as $1/T$, and the coefficients of those polynomials grow geometrically in $T$, so by part (i) accuracy improves only logarithmically in $N$: the rate $1/\log N$ of Theorem~\textup{6.7 in Paper~1}. \hfill\ProvedIn{PAPER 1}
\end{enumerate}
\end{corollary}

The link in (e) between coefficient growth and the logarithmic rate is a heuristic reading of Paper~1's proof, not a separate result; making it a corollary of part (iv) is part of the target.

\subsection{Partial lower bound and the gap in part (iv)}\label{sec:dich_partial}

A repaired version of part (iv) remains a \Target{}. The original uniform statement is refuted below. The following argument establishes only a parametric obstruction, with a constant depending on the functional. It does not establish the coefficient-energy factor, and has not been formalized in Lean.

\paragraph{A fixed-atom lower bound.}
Suppose $B>0$, $\sigma>0$, $\alpha<1/4$, and choose $r_*\in[0,\bar r]$ with $g(r_*)\ne0$. For sufficiently small $\epsilon$, any valid certificate must satisfy
\[
N_{\rm cert}(\epsilon;F_g,\mathcal M(\bar r,B))>
\frac{\sigma^2(1-r_*^2)|g(r_*)|^2}{9\epsilon^2}.
\]
Indeed, put $\lambda=3\epsilon/|g(r_*)|$ and require $\lambda\le B/2$. The measures
$M_0=(B/4)\delta_{r_*}$ and $M_1=(B/4+\lambda)\delta_{r_*}$ are admissible and their functional values differ by $3\epsilon$. For any observed depth $L$, their Gaussian experiment divergence obeys
\[
D_{\rm KL}(P_0^N\Vert P_1^N)
=\frac{N\lambda^2}{2\sigma^2}\sum_{t<L}r_*^{2t}
\le\frac{N\lambda^2}{2\sigma^2(1-r_*^2)}.
\]
If $N$ is no greater than the displayed lower bound, the divergence is at most $1/2$, and Pinsker's inequality gives total variation at most $1/2$. An interval of half-width at most $\epsilon$ cannot contain both functional values. Coverage at least $1-\alpha$ under both measures would therefore imply total variation at least $1-2\alpha>1/2$, a contradiction. The divergence bound is uniform in $L$.

\paragraph{Why the proposed matching proof does not close the gap.}
Taylor truncation gives an upper bound on best polynomial approximation error, $B E_T(g)\le\tau_g(T)$; it gives no reverse bound. Thus part (ii) does not show that every procedure needs Taylor depth $T_g(\epsilon)$. Furthermore, a Gauss/Gauss--Radau pair with $T$ nodes matches all moments through degree $2T-2$. For $t<T$, its tap difference is exactly zero, rather than a nonzero vector aligned with $(a_0,\ldots,a_{T-1})$. Positivity of the coefficients does not supply the missing construction. Finally, when $L>T$, a valid testing lower bound must control divergence from all observed taps, not just a selected prefix. The fixed-atom argument above does so, but does not yield $A_g(T_g(\epsilon))$.

\paragraph{Counterexample to the uniform statement as written.}
The constants in part (iv) are stated to depend only on $(\bar r,B,\alpha)$, uniformly over the regular coefficient class. This uniform claim fails. Fix $\bar r>0$, $B>0$, an accuracy $\epsilon_0>0$, and $\beta\in(0,1)$, and choose $\eta>0$ with $\eta/(1-\beta\bar r)<1/2$. For any positive integer $m$, define
\[
g_m(r)=\frac{\epsilon_0}{B}\left[(1-r/\bar r)^m+\frac{\eta}{1-\beta r}\right].
\]
Its coefficients are strictly positive beyond degree $m$ and have at most polynomial growth. On the support interval, $g_m$ is non-negative and bounded above by $(\epsilon_0/B)(1+\eta/(1-\beta\bar r))$. Thus every $F_{g_m}(M)$ lies in a known interval of length less than $3\epsilon_0/2$. That interval certifies the functional without any data, so $N_{\rm cert}\le1$ even if the definition requires a positive number of experiments.

Yet $T_{g_m}(\epsilon_0)=m+1$. For every $1\le T\le m$, evaluating the polynomial tail at $r=\bar r$ gives
\[
\sum_{t=T}^m(-1)^t\binom mt=(-1)^T\binom{m-1}{T-1}.
\]
After adding the discounted tail, the reverse triangle inequality yields $\tau_{g_m}(T)\ge\epsilon_0(1-\eta/(1-\beta\bar r))>\epsilon_0/2$. At $T=m+1$ only the discounted tail remains and is below $\epsilon_0/2$. For any fixed $c'>0$, eventually $c'(m+1)>1$, so the coefficient energy includes the coefficient of $r$ and
\[
A_{g_m}(c'T_{g_m}(\epsilon_0))\epsilon_0^{-2}
\ge B^{-2}(m/\bar r-\eta\beta)^2\longrightarrow\infty.
\]
For fixed $\sigma>0$ this contradicts any positive uniform constant $c$ in part (iv). The accuracy $\epsilon_0$ can be chosen arbitrarily small. This refutes the stated bound over the full regular class at a common accuracy; it does not refute an asymptotic statement with a functional-dependent accuracy threshold, nor the restricted all-non-negative coefficient case. A repaired statement must make its additional normalization or its functional dependence explicit. The original claim is retained visibly as refuted; a repaired lower-bound target must state its additional assumptions explicitly.

\paragraph{Deterministic checks and remaining target.}
The script \texttt{checks/check\_dich\_lower.py} records valid fixed-atom pairs for the safety, single-tap and discounted functionals, their functional gaps and divergences, and the exactness residuals of quadrature pairs in a JSON result file in the same directory. It checks the parametric scaling of the explicit testing budget. This scaling is a consequence of the construction, not independent evidence for a minimax matching rate. The script also evaluates the cancellation family above. No empirical or simulation test is performed. A reformulation of the uniform bound and the non-negative coefficient case remain open, as does the proposed deduction for Corollary~\ref{cor:cases}(e).

\section{Disclosure changes the contract, not the question}\label{sec:fw_disclosure}

Disclosure regimes are observation contracts. Class-tagged flow reveals each class's own moments, which turns identification into an ordinary estimation problem and is worth a factor $(\Delta\eta)^{-(4K-4)}$ in samples (Theorem~\textup{8.2 in Paper~1}). For safety the value of disclosure is different: it raises safe capacity exactly when it rules out the extremal compatible economy, and it is worth most when memory is long, where Theorem~\ref{thm:dich}(iii) says public data certify nothing (Theorem~\textup{7.1 in Paper~2}). Supervision that wants safety should therefore ask for the disclosures that bound memory or reveal $\Phi$, not for identities.

\section{A robust safety bound: nonlinear and changing agents}\label{sec:bl_robust}

Papers~1 and~2 work with agents whose rules are linear filters with fixed memory and weight. Real agents cap positions, trade only beyond thresholds, adapt their gains, and keep learning while the operator watches. This section and the next show that the safety side survives all of this under one condition: a bound on the agents' total sensitivity.


Return to the economy of Definition~\ref{def:cohort} and let everything about the agents vary over time.

\begin{assumption}[Bounded total sensitivity]\label{ass:sector}
For each cohort $i$ and period $t$, the memory rate $r_i(t)\in[0,1)$ is arbitrary, the weight satisfies $|q_i(t)|\le\bar q_i$, and the response rule satisfies the sector condition $|\varphi_{i,t}(y)|\le s_i|y|$ for all $y$. Write $\bar\Phi=\sum_i\bar q_i s_i$ for the total sensitivity.
\end{assumption}

The sector condition covers saturating momentum ($s_i=1$ for $\tanh$), position limits, threshold and dead-zone rules ($y\mathbf 1\{|y|>\tau\}$), and any rule whose response never exceeds a multiple of its signal. It does not require the rule to be smooth, monotone, stationary or known.

\begin{theorem}[Robust safety bound]\label{thm:sector}
Under Assumption~\ref{ass:sector}, if $|a|+\kappa\bar\Phi<1$ and the economy starts at rest, then for every shock sequence with $|w_t|\le\bar w$,
\[
\sup_{t\ge0}|x_t|\;\le\;\bar w\,\frac{\kappa}{1-|a|-\kappa\bar\Phi}.
\]
When $a\ge0$, this is the worst-case loss bound of the linear stationary economy with weights $\bar q_is_i$, so the summed-gain guard of Paper~2 computed at total sensitivity $\bar\Phi$ is safe for every economy satisfying the assumption. \hfill\Proved
\end{theorem}

\paragraph{Proof.} See Appendix~\ref{app:paper3proof2}.


\paragraph{What the theorem says.} An operator who wants a safe admission level does not need to know how many agents there are, what their memories are, whether their rules are linear, or whether they change over time. It needs one number, an upper bound on the agents' total sensitivity $\bar\Phi$. The bound is tight: when the agents are linear, stationary, positive and at their maximal weights, it is attained. Nonlinearity, memory changes and negative feedback can only make the market safer than the bound says.

\begin{remark}[What the theorem does not cover]
Rules that amplify their signal, $|\varphi(y)|>s|y|$ for the declared $s$, break the bound; so do changes in the market's own persistence $a$ or impact $\kappa$, which the operator must certify separately (Paper~2, Section~\textup{5 in Paper~2}). The bound controls the peak of the price path, which is what a fixed position's loss depends on; dynamic positions are handled by the balance-sheet guard of Paper~2.
\end{remark}

\begin{proposition}[Several assets]\label{prop:multiasset}
Let $x_t\in\R^d$ evolve as $x_{t+1}=Ax_t+\mathcal K\bigl(\sum_iQ_i(t)\varphi_{i,t}(y_{i,t})+w_{t+1}\bigr)$, with memory states $y_{i,t}=R_i(t)y_{i,t-1}+(I-R_i(t))x_t$, $R_i(t)$ diagonal with entries in $[0,1)$, sector conditions $|\varphi_{i,t}(y)|\le s_i|y|$ componentwise, and $\|Q_i(t)\|_\infty\le\bar q_i$. If $\|A\|_\infty+\|\mathcal K\|_\infty\sum_i\bar q_is_i<1$, then $\sup_t\|x_t\|_\infty\le\bar w\,\|\mathcal K\|_\infty/(1-\|A\|_\infty-\|\mathcal K\|_\infty\sum_i\bar q_is_i)$ for $\|w_t\|_\infty\le\bar w$. \hfill\Proved
\end{proposition}

\paragraph{Proof.} See Appendix~\ref{app:paper3proof3}.


The induced-norm bound discards the pattern of cross-impact. For positive cross-impact, the following comparison gives the sharper replacement. Its proof is written here; full Lean formalization remains outstanding, so its target status is retained.

\begin{proposition}[Positive cross-impact comparison]\label{prop:multiasset_positive}
Assume the economy starts at rest, $A,\mathcal K,\overline Q_i$ are entrywise non-negative, $|Q_i(t)|\le\overline Q_i$ entrywise, and $|\varphi_{i,t}(y)|\le S_i|y|$ for a non-negative diagonal matrix $S_i$. Memory matrices $R_i(t)$ are diagonal with entries in $[0,1)$. Put $G=A+\mathcal K\sum_i\overline Q_i S_i$. If $\rho(G)<1$, then for every $\|w_t\|_\infty\le\bar w$,
\[
\sup_t\|x_t\|_\infty\le
\bar w\,\bigl\|(I-G)^{-1}\mathcal K\bigr\|_\infty.
\]
This bound is sharp over the stated class. \hfill\Target
\end{proposition}

\paragraph{Proof.} See Appendix~\ref{app:paper3proof4}.


\paragraph{Registered numerical validation.}
% Generated by checks/export_multiasset_numbers.py from R_multiasset_recovery.json
\newcommand{\MultiassetCases}{128}
\newcommand{\MultiassetPassed}{128}

After the original registered case archive was overwritten by a concurrent writer, a separate follow-up specification and hash were pushed before a new seeded set of instances was generated. The original test was not reopened, and the replacement external results were not used. The follow-up unlock counter was advanced before its single opening. The joint primary required both componentwise enclosure for nonlinear, time-varying cohorts and convergence to the bound in the maximal positive linear instance. It passed in \MultiassetPassed{} of \MultiassetCases{} registered random instances. The generated matrices, tolerances and diagnostics are recorded in \nolinkurl{checks/R_multiasset_recovery.json}. This is numerical evidence under the registered construction; it does not prove the result for every matrix or replace the outstanding Lean formalization of the positive resolvent and sharpness steps.

\section{Certificates expire: learning agents}\label{sec:bl_learning}

Theorem~\ref{thm:sector} allows weights to change, provided they stay below their bounds. When agents learn, the bounds themselves are the issue: a certificate issued today says nothing about weights next month. The theorem turns this into a computable lifetime.

\begin{corollary}[Certificate lifetime]\label{thm:lifetime}
Suppose a certificate at time $0$ establishes $|q_i(0)|\le q_i^+$ for each cohort, and the guard's loss budget is met for every total sensitivity up to $\Phi_{\max}$, that is $\kappa\bar w/(1-a-\kappa\Phi_{\max})$ times the position equals the budget. If each weight changes by at most $\gamma_i$ per period, $|q_i(t+1)-q_i(t)|\le\gamma_i$, and $\sum_iq_i^+s_i\le\Phi_{\max}$ with $\sum_i\gamma_is_i>0$, then the guard remains safe for at least
\[
T_{\rm valid}=\Bigl\lfloor\frac{\Phi_{\max}-\sum_iq_i^+s_i}{\sum_i\gamma_is_i}\Bigr\rfloor
\]
periods under the sector bounds. \hfill\Proved
\end{corollary}

\paragraph{Proof.} See Appendix~\ref{app:paper3proof5}.


The lifetime is the certificate's safety margin divided by the agents' speed of learning. This is a guaranteed lower bound, not a claim that loss first crosses the budget at this time. Realized safety can last much longer if learning slows, changes sign, or stops below the sensitivity budget. When $\sum_i\gamma_i s_i=0$, the sensitivity envelope does not grow and the same certificate remains valid for every horizon under the stated sector and plant assumptions; the displayed quotient is not used. If the initial sensitivity bound already exceeds $\Phi_{\max}$, this argument supplies no valid initial certificate. Recertification frequency can therefore be based on the certified sensitivity margin and the learning-speed bound, while additional information about the realized learning path may justify a longer horizon.

\paragraph{Formal scope.}
The finite-cohort, finite-horizon safety induction is checked by \nolinkurl{CurrentMM.thesis_sector_safety}. The drift, floor and sensitivity calculations are checked by \nolinkurl{CurrentMM.thesis_lifetime_envelope} and its helper lemmas. The deterministic certification arithmetic is checked by \nolinkurl{CurrentMM.thesis_variance_radius} and \nolinkurl{CurrentMM.thesis_interval_half_width}; the general finite-measure bias step is checked by \nolinkurl{CurrentMM.thesis_general_measure_bias}, under explicit integrability and almost-everywhere tail hypotheses. These declarations do not formalize Gaussian coverage, the derivation of analytic tail or integrability hypotheses, the matching lower bound, or the sharp positive cross-impact proposition. The coverage register records these boundaries separately.

\paragraph{Identification under learning, by contrast.} Paper~1, Section~\textup{7 in Paper~1}, shows that trained online learners obey the same $\delta^{-4}$ law for separating memory rates when the competing populations' weights match, and that populations whose learned weights differ are told apart by those weights instead. Learning therefore leaves identification as hard as before while making certification only perishable, not harder per renewal.

\section{Markets of trading algorithms that learn}\label{sec:bl_ai}

Markets in which learning algorithms, including reinforcement-learning and language-model agents, trade with each other are where the results of this chapter matter most. Two of them apply directly. Certification of safety needs only a bound on total sensitivity (Theorem~\ref{thm:sector}); and the certificate lasts for its safety margin divided by the agents' learning speed (Corollary~\ref{thm:lifetime}). The second result gives a supervisory rule with no counterpart in identity-based regimes: recertify at a frequency set by how fast the agents update.

\pendbox{multi-agent simulation}{Simulate markets of reinforcement-learning trading agents in the price model of Paper~2, measure the realized lifetime of certificates as agents keep learning, and compare it with Corollary~\ref{thm:lifetime}. The specification (agents, learning rates, horizons and pass criteria) is to be registered before the simulation is run.}

\section{Agents who respond to the guard}\label{sec:bl_strategic}

A guard changes the environment it certifies. If agents learn that the operator throttles admission when measured sensitivity is high, they may change how they trade.

\begin{conjecture}[Guard-consistent equilibrium]\label{conj:strategic}
Suppose each cohort chooses its weight $q_i\in[0,\bar q_i]$ to maximize expected trading profit net of a convex cost, given an admission rule that throttles the dealer's position when certified total sensitivity exceeds $\Phi_{\max}$, and given the equilibrium reversion supplied by dealers (Proposition~\textup{8.1 in Paper~2}). Then (i) an equilibrium exists; (ii) in every equilibrium the guard is safe; and (iii) equilibrium total sensitivity is weakly lower than without the guard. \hfill\Refuted
\end{conjecture}

\begin{remark}[Counterexample to existence under a discontinuous guard]\label{rem:strategic_counterexample}
Part (i) of the conjecture is false under the stated assumptions. Consider one cohort with $q\in[0,1]$, a perfect certificate $\Phi=q$, and admitted dealer position $I(q)=1$ for $q\le1/2$ and $I(q)=0$ for $q>1/2$. This is an admission rule that throttles above the certified threshold. Let expected trading profit, as a function of weight and admitted position, be
\[
\pi(q,I)=(1-I)(1-q),\qquad c(q)=0.
\]
The profit primitive is continuous in $(q,I)$ and the zero cost is convex. After composition with the guard, utility is zero at $q\le1/2$ and equals $1-q$ at $q>1/2$. If $q\le1/2$, the admissible deviation $q'=3/4$ strictly increases utility. If $q>1/2$, the admissible deviation $q'=(q+1/2)/2$ lies strictly between $1/2$ and $q$ and also strictly increases utility. Thus no weight maximizes utility: the one-player game has no pure equilibrium despite its compact strategy set, convex cost, and exact sensitivity certificate.

Randomization does not restore existence. Utility is strictly below $1/2$ everywhere, so its expectation under any probability distribution is strictly below $1/2$; a pure weight sufficiently close to the threshold from above earns more. The universal profitable-deviation statement and nonexistence of a pure maximizer are formalized as \nolinkurl{CurrentMM.thesis_strategic_profitable_deviation} and \nolinkurl{CurrentMM.thesis_strategic_no_equilibrium}. The probability-integral step for arbitrary mixed strategies is a written argument, not part of those Lean declarations. Exact rational diagnostics are recorded separately in \texttt{checks/R\_strategic.json}.

This is a counterexample at the level of the payoff assumptions currently stated: profits are permitted to increase when dealer admission falls. It does not refute existence in a separately specified market P\&L model with additional payoff restrictions. The failure concerns equilibrium existence, not the safety invariant. For example take $a=0$, $\kappa=1/4$, and $\bar w=1$. On admitted profiles $q\le1/2$, unit position loss is bounded by $2/7$; on throttled profiles the position is zero. Every admitted decision therefore respects this budget, although the agents' game has no equilibrium. A revised strategic theorem needs assumptions that rule out the best-response discontinuity. Convex cost alone also supplies no ordering of equilibrium sensitivity with and without a guard; that comparative-static claim remains open under suitable additional payoff restrictions.
\end{remark}

\section{What the framework does not cover}\label{sec:fw_scope}

The framework is linear in the observable layer and assumes the reset contract for its sharpest statements. Sections~\ref{sec:bl_robust} and~\ref{sec:bl_learning} extend the safety results to sector-bounded nonlinear and to learning agents; strategic equilibrium existence fails without additional assumptions (Remark~\ref{rem:strategic_counterexample}); existence for suitably regularized guards remains open. Passive observational data are slower but obey the same polynomial order for identification (Theorem~\textup{6.9 in Paper~1}); a finite-sample theory of certification from observational order-book data, with dependence and non-stationarity, is not developed here, and the empirical certificates of Paper~2 are bootstrap diagnostics rather than guarantees.

\appendix
\section{Proofs}
\subsection{Proof for result~\ref{thm:dich}}\label{app:paper3proof1}
\begin{proof}[Proof of (i)--(iii)]
(i) Let $\bar k_t$ be the average of the $N$ replicates of tap $t$ and $T=T_g(\epsilon)$. The estimator $\hat F=\sum_{t<T}a_t\bar k_t$ has mean $\sum_{t<T}a_tk_t$ and variance $\sigma^2A_g(T)/N$. Its bias is $|\sum_{t\ge T}a_tk_t|=|\int\sum_{t\ge T}a_tr^t\,dM(r)|\le\tau_g(T)\le\epsilon/2$ for every $M\in\mathcal M(\bar r,B)$. The interval $\hat F\pm(\epsilon/2+z\sigma\sqrt{A_g(T)/N})$ therefore covers $F_g(M)$ with probability at least $1-\alpha$, and its half-width is at most $\epsilon$ once $z\sigma\sqrt{A_g(T)/N}\le\epsilon/2$.

(ii) If $E_T(g)=0$, the asserted zero lower bound is immediate. Otherwise, by duality for best uniform approximation, $E_T(g)=\max\{\int g\,d\nu\}$ over signed measures $\nu$ on $[0,\bar r]$ with total variation at most one that annihilate every polynomial of degree below $T$. Take a maximizer and write $\nu=\nu^+-\nu^-$. Orthogonality to constants gives $\nu^+([0,\bar r])=\nu^-([0,\bar r])=:c\le1/2$. Set $M=(B/c)\nu^+$ and $M'=(B/c)\nu^-$: both lie in $\mathcal M(\bar r,B)$, their first $T$ moments coincide, and $F_g(M)-F_g(M')=(B/c)E_T(g)\ge2B\,E_T(g)$. Two measures that produce identical data cannot both be covered by an interval of half-width below half their distance.

(iii) Compare $M_0=0$ with $M_1=c\,\delta_r$. Under $M_1$ the taps are $k_t=c\,r^t$, so the Kullback--Leibler divergence between the two data distributions over any number of taps is at most $Nc^2/(2\sigma^2(1-r^2))$. Choose $c=\sigma\sqrt{(1-r^2)/N}$, which is at most $B$ for $r$ near one; the divergence is then at most $1/2$, so by Pinsker's inequality the total variation distance between the two data distributions is at most $1/2$. The functional differs by $c\,|g(r)|=\sigma\sqrt{(1-r^2)/N}\,|g(r)|\ge\sigma\sqrt{(1-r)/N}\,|g(r)|$, which is unbounded along a sequence $r\uparrow1$ by assumption. Suppose an interval of half-width $\epsilon$ covered $F_g$ with probability at least $1-\alpha$ under both economies. Once the two values differ by more than $2\epsilon$, the event that the interval contains $F_g(M_0)$ has probability at least $1-\alpha$ under $M_0$ and at most $\alpha$ under $M_1$, so the total variation distance is at least $1-2\alpha>1/2$, a contradiction.
\end{proof}
\subsection{Proof for result~\ref{thm:sector}}\label{app:paper3proof2}
\begin{proof}
Let $Z_t=\max_{s\le t}|x_s|$. Each memory state is a convex combination of its previous value and the current state, so $|y_{i,t}|\le r_i(t)|y_{i,t-1}|+(1-r_i(t))|x_t|\le\max(|y_{i,t-1}|,|x_t|)$, and by induction from $y_{i,0}=0$, $|y_{i,t}|\le Z_t$. Then
\[
|x_{t+1}|\le|a|\,|x_t|+\kappa\Bigl(\sum_i|q_i(t)|\,|\varphi_{i,t}(y_{i,t})|+|w_{t+1}|\Bigr)\le(|a|+\kappa\bar\Phi)\,Z_t+\kappa\bar w .
\]
Let $Z^*=\kappa\bar w/(1-|a|-\kappa\bar\Phi)$. If $Z_t\le Z^*$ then $|x_{t+1}|\le(|a|+\kappa\bar\Phi)Z^*+\kappa\bar w=Z^*$, so $Z_{t+1}\le Z^*$; and $Z_0=0$. For the second claim, with $a\ge0$ the linear stationary economy with weights $\bar q_is_i$ is a positive system, its impulse response is non-negative, and its worst-case peak response to shocks bounded by $\bar w$ is $\bar w\sum_th_t=\bar w\,\kappa/(1-a-\kappa\bar\Phi)$ (equation~\eqref{eq:fw_S}). The bound of the theorem therefore coincides with the quantity the guard controls (Lemma~\textup{6.1 in Paper~2}).
\end{proof}
\subsection{Proof for result~\ref{prop:multiasset}}\label{app:paper3proof3}
\begin{proof}
The proof of Theorem~\ref{thm:sector} goes through with absolute values replaced by the sup norm and scalars by induced sup norms: each component of $y_{i,t}$ is a convex combination of its past value and the corresponding component of $x_t$.
\end{proof}
\subsection{Proof for result~\ref{prop:multiasset_positive}}\label{app:paper3proof4}
\begin{proof}
The Neumann series $\sum_{j\ge0}G^j=(I-G)^{-1}$ has non-negative entries. Write $\mathbf1$ for the vector of ones and set $z=\bar w(I-G)^{-1}\mathcal K\mathbf1\ge0$. Then $Gz+\bar w\mathcal K\mathbf1=z$. Suppose $|x_t|\le z$ and $|y_{i,t}|\le z$ componentwise. The plant and the sector bounds give
\[
|x_{t+1}|\le A z+\mathcal K\Bigl(\sum_i\overline Q_i S_i z+\bar w\mathbf1\Bigr)=z.
\]
Each next memory coordinate is a convex combination of its previous value and the corresponding next state coordinate, so its absolute value is also bounded by the corresponding coordinate of $z$. Induction from rest proves the componentwise bound for all times. Because the resolvent times $\mathcal K$ is non-negative, its induced sup norm is its maximum row sum, and $\|z\|_\infty$ equals the displayed bound.

For sharpness, choose $R_i(t)=0$, $Q_i(t)=\overline Q_i$, $\varphi_{i,t}(y)=S_i y$, and $w_t=\bar w\mathbf1$. The state then satisfies $x_{t+1}=Gx_t+\bar w\mathcal K\mathbf1$, so from rest $x_t=\bar w\sum_{j<t}G^j\mathcal K\mathbf1$ increases componentwise to $z$. Its supremum norm therefore approaches the bound. Sharpness means a supremum over the admitted model class, not attainment at a finite time in every economy.
\end{proof}
\subsection{Proof for result~\ref{thm:lifetime}}\label{app:paper3proof5}
\begin{proof}
On $[0,T_{\rm valid}]$, $|q_i(t)|\le q_i^++\gamma_iT_{\rm valid}=:\bar q_i$ and $\sum_i\bar q_is_i\le\Phi_{\max}$. Theorem~\ref{thm:sector} applies on that horizon, since its induction uses only periods up to the current one.
\end{proof}
\begin{thebibliography}{99}\setlength{\itemsep}{2pt}\small
\bibitem[Brunnermeier and Pedersen(2009)]{t:brunnermeier2009} Brunnermeier, M.~K. and Pedersen, L.~H. (2009). Market liquidity and funding liquidity. \emph{Review of Financial Studies} 22(6), 2201--2238.
\bibitem[Cai and Low(2011)]{t:cai2011} Cai, T.~T. and Low, M.~G. (2011). Testing composite hypotheses, Hermite polynomials and optimal estimation of a nonsmooth functional. \emph{Annals of Statistics} 39(2), 1012--1041.
\bibitem[Cand\`es and Fernandez-Granda(2014)]{t:candes2014} Cand\`es, E.~J. and Fernandez-Granda, C. (2014). Towards a mathematical theory of super-resolution. \emph{Communications on Pure and Applied Mathematics} 67(6), 906--956.
\bibitem[Farina and Rinaldi(2000)]{t:farina2000} Farina, L. and Rinaldi, S. (2000). \emph{Positive Linear Systems: Theory and Applications}. Wiley.
\bibitem[Heinrich and Kahn(2018)]{t:heinrich2018} Heinrich, P. and Kahn, J. (2018). Strong identifiability and optimal minimax rates for finite mixture estimation. \emph{Annals of Statistics} 46(6A), 2844--2870.
\bibitem[Kirilenko et~al.(2017)]{t:kirilenko2017} Kirilenko, A., Kyle, A.~S., Samadi, M. and Tuzun, T. (2017). The Flash Crash: High-frequency trading in an electronic market. \emph{Journal of Finance} 72(3), 967--998.
\bibitem[Lepski et~al.(1999)]{t:lepski1999} Lepski, O., Nemirovski, A. and Spokoiny, V. (1999). On estimation of the $L_r$ norm of a regression function. \emph{Probability Theory and Related Fields} 113(2), 221--253.
\bibitem[Ljung(1999)]{t:ljung1999} Ljung, L. (1999). \emph{System Identification: Theory for the User}, 2nd edition. Prentice Hall.
\bibitem[Manski(2003)]{t:manski2003} Manski, C.~F. (2003). \emph{Partial Identification of Probability Distributions}. Springer.
\bibitem[Megretski and Rantzer(1997)]{t:megretski1997} Megretski, A. and Rantzer, A. (1997). System analysis via integral quadratic constraints. \emph{IEEE Transactions on Automatic Control} 42(6), 819--830.
\bibitem[Moitra(2015)]{t:moitra2015} Moitra, A. (2015). Super-resolution, extremal functions and the condition number of Vandermonde matrices. In \emph{Proceedings of the 47th Annual ACM Symposium on Theory of Computing}, 821--830.
\bibitem[Rantzer(2015)]{t:rantzer2015} Rantzer, A. (2015). Scalable control of positive systems. \emph{European Journal of Control} 24, 72--80.
\bibitem[Wu and Yang(2020)]{t:wu2020} Wu, Y. and Yang, P. (2020). Optimal estimation of Gaussian mixtures via denoised method of moments. \emph{Annals of Statistics} 48(4), 1981--2007.

\bibitem[Niroula and Niroula(2026)]{p3:paper1} Niroula, P. and Niroula, P. (2026). \emph{Auditable Memory and Financial Identification Limits in Electronic Limit Order Books}. Companion working paper.
\bibitem[Niroula et~al.(2026)]{p3:paper2} Niroula, P., Alkhushayni, S., Chen, Y., Bhavsar, M. and Niroula, P. (2026). \emph{From Human Markets to Machine Markets: Transient Amplification, an Exact Depth Condition, and a Maximal Pre-Trade Guard}. Companion working paper.
\end{thebibliography}

\end{document}
